\documentclass[11pt]{article}
\usepackage[margin=1in]{geometry}
\usepackage[T1]{fontenc}
\usepackage{lmodern}
\usepackage{amsmath,amssymb,amsthm,mathrsfs}
\usepackage{graphicx}
\usepackage[hidelinks]{hyperref}

\newtheorem{theorem}{Theorem}
\newtheorem{lemma}[theorem]{Lemma}
\newtheorem{corollary}[theorem]{Corollary}
\newcommand{\R}{\mathbb R}
\newcommand{\Pj}{\mathbb{RP}^{1}}
\newcommand{\Crit}{\operatorname{Crit}}
\newcommand{\rank}{\operatorname{rank}}
\newcommand{\Leb}{\mathcal L^2}

\title{A Global Quadratic Sard Theorem\\
\large for maps from a Hilbert space to the plane}
\author{Automated mathematical discovery run \texttt{fa\_banach\_001}\\
Lane 6, \texttt{agent\_lane\_06}}
\date{9 August 2026}

\begin{document}
\maketitle

\noindent\textbf{Status.}
Candidate partial result, likely valid, subject to human review.

\medskip
\noindent\textbf{Source question.}
Lerario, Rizzi, and Tiberio define a class
$\mathscr P_d^m(H)$ of polynomial maps from a real Hilbert space to $\R^m$,
with a weakly continuous and locally Lipschitz differential. After constructing
counterexamples in degrees at least three, their Remark~3 says:
\begin{quote}
The case $d=2$ is special. The authors expect every
$f\in\mathscr P_2^m(H)$ to satisfy the Sard property globally on $H$.
\end{quote}
The exact source passage is on page 3 of arXiv:2407.02296v4 and is reproduced
below.

\begin{center}
\includegraphics[width=0.94\textwidth]{figures/open_problem_crop.png}
\end{center}

The result here proves the expectation completely for codomain dimension two,
with no commutativity or regularity assumption on the self-adjoint Hessian
pencil.

\section{Statement}

Every $f\in\mathscr P_2^2(H)$ has a representation
\begin{equation}\label{eq:representation}
 f_j(x)=c_j+2\langle a_j,x\rangle+\langle A_jx,x\rangle,
 \qquad j=1,2,
\end{equation}
where $A_1,A_2$ are compact self-adjoint operators. Put
\[
 K=\ker A_1\cap\ker A_2,
 \qquad
 L_Kz=\bigl(\langle a_1,z\rangle,\langle a_2,z\rangle\bigr).
\]
\begin{theorem}[global quadratic Sard theorem in target dimension two]
\label{thm:main}
Let $H$ be a real Hilbert space. Every $f\in\mathscr P_2^2(H)$ satisfies the
global Sard property:
\[
 \Leb\bigl(f(\Crit(f))\bigr)=0.
\]
\end{theorem}

\section{Proof intuition}

A critical point has a projective multiplier $[\lambda_1:\lambda_2]$ for
which the scalar quadratic $\lambda_1f_1+\lambda_2f_2$ is stationary. On a
bounded ball, all pairs consisting of a stationary point and its multiplier
form a compact incidence set when the ball is given its weak topology.

The key observation compares two such pairs. If $x,z$ have nearby
multipliers $\lambda,\mu$, the two exact quadratic Taylor formulas imply
\[
 (\lambda+\mu)\cdot(f(z)-f(x))
 =\langle A_{\lambda-\mu}(z-x),z-x\rangle.
\]
Consequently their critical values lie in a thin strip: its normal thickness
is bounded by a constant times the distance between the multipliers.

The pencil can be singular in uncountably many directions, so uniqueness of
the stationary point is false in general. The decisive replacement is that
only countably many multiplier directions can have a stationary fiber on
which $f$ is nonconstant. Indeed, on an affine chart of the multiplier circle,
write $B_t=A_2-tA_1$ and $K_t=\ker B_t$. If the quadratic part varies on
$K_t$, vectors from distinct $K_t$ give a biorthogonal system. If only the
linear part varies, the stationary gradients give another biorthogonal
system. Every biorthogonal system in a separable Hilbert space is countable.

Away from those active directions, a multiplier has a unique critical
\emph{value}, even when it has many stationary points. Compactness of the
incidence set makes that value uniformly continuous on compact multiplier
sets. Splitting the multiplier circle into arcs of length $\delta$ covers the
image by $O(\delta^{-1})$ rectangles, each with one side $O(\delta)$ and the
other $o(1)$. Countably many active directions can be put inside open arcs of
arbitrarily small total length. The strip estimate finishes the proof.

\section{Proof}

\subsection{Quadratic representation and compact Hessians}

For a scalar component of $f$, subtract its value and differential at the
origin. Its restriction to every finite-dimensional subspace is a homogeneous
quadratic polynomial. Polarization therefore gives a continuous symmetric
bilinear form on $H$, represented by a bounded self-adjoint operator. This
gives \eqref{eq:representation}. Moreover,
\[
 Df_j(x)h=2\langle A_jx+a_j,h\rangle.
\]
If $x_n\rightharpoonup0$, weak continuity of $Df$ into the operator norm gives
$\|A_jx_n\|\to0$. The weak-null sequence characterization of compact
operators on Hilbert space shows that $A_j$ is compact.

Only a separable part of $H$ matters. Indeed, the closed span of
$a_1,a_2,\operatorname{ran}A_1,\operatorname{ran}A_2$ is separable, and $f$ is
constant on its orthogonal complement. We may consequently assume $H$ is
separable. On bounded subsets, each compact quadratic form
$x\mapsto\langle A_jx,x\rangle$ is weakly continuous. Thus $f$ is weakly
continuous on bounded balls.

\subsection{The common kernel}

On $K$, the restriction of every derivative is the fixed linear map $2L_K$.
If $\rank L_K=2$, every $Df(x)$ is onto and there are no critical points.

Suppose $\rank L_K=1$. Any multiplier $[\lambda]$ at a critical point must
annihilate $L_K(K)$, so the projective multiplier is unique. Write
\[
 A_\lambda=\lambda_1A_1+\lambda_2A_2,
 \qquad a_\lambda=\lambda_1a_1+\lambda_2a_2.
\]
If $x,z$ are critical, both solve $A_\lambda u+a_\lambda=0$. Hence
$A_\lambda(z-x)=0$, and the exact Taylor formula gives
\[
 \lambda\cdot(f(z)-f(x))
 =\langle A_\lambda(z-x),z-x\rangle=0.
\]
All critical values lie on one affine line and are null.

It remains to consider $L_K=0$. The linear and quadratic terms vanish on $K$,
including all cross terms, so $f$ factors exactly through $K^\perp$. We replace
$H$ by $K^\perp$; henceforth $K=\{0\}$.

\subsection{The multiplier incidence and the strip estimate}

Let $P=\Pj$ be the projective multiplier circle. For $R>0$ define
\[
 \mathcal I_R=
 \{([\lambda],x)\in P\times B_H(R):A_\lambda x+a_\lambda=0\}.
\]
The equation is invariant under changing the sign of a representative.
Give $B_H(R)$ its weak topology. It is compact and metrizable. If
$[\lambda_n]\to[\lambda]$ and $x_n\rightharpoonup x$, choose locally aligned
unit representatives. Compactness of $A_\lambda$ and norm continuity of the
pencil give
\[
 A_{\lambda_n}x_n+a_{\lambda_n}\longrightarrow A_\lambda x+a_\lambda
 \quad\hbox{in norm}.
\]
Thus $\mathcal I_R$ is compact. Its image under $([\lambda],x)\mapsto f(x)$
is exactly $f(\Crit(f)\cap B_H(R))$.

We need one quantitative identity. Take $([\lambda],x),([\mu],z)$ in
$\mathcal I_R$ and choose aligned unit representatives. Put $h=z-x$. Taylor
expansion at the two stationary points gives
\begin{align*}
 \lambda\cdot(f(z)-f(x))&=\langle A_\lambda h,h\rangle,\\
 \mu\cdot(f(z)-f(x))&=-\langle A_\mu h,h\rangle.
\end{align*}
Adding them yields the exact identity
\begin{equation}\label{eq:slabidentity}
 (\lambda+\mu)\cdot(f(z)-f(x))
 =\langle A_{\lambda-\mu}h,h\rangle.
\end{equation}
Since $\|h\|\leq2R$ and $f(B_H(R))$ is bounded, \eqref{eq:slabidentity}
implies
\begin{equation}\label{eq:slab}
 |\lambda\cdot(f(z)-f(x))|
 \leq C_R\|\lambda-\mu\|
\end{equation}
for a constant depending only on $R,f$. Thus values whose multipliers lie in
one short projective arc of length $\ell$ lie in a strip of normal width at
most $2C_R\ell$.

\subsection{Only countably many stationary fibers are active}

We use the following elementary separability fact.

\begin{lemma}[biorthogonal families]\label{lem:biorthogonal}
Let $H$ be a separable Hilbert space. If
$(u_\alpha,v_\alpha)_{\alpha\in I}\subset H\times H$ satisfies
\[
 \langle u_\alpha,v_\beta\rangle=
 \begin{cases}
 1,&\alpha=\beta,\\
 0,&\alpha\ne\beta,
 \end{cases}
\]
then $I$ is countable.
\end{lemma}

\begin{proof}
Rescale so that $\|v_\alpha\|=1$. For $n\geq1$, the vectors $v_\alpha$ with
$\|u_\alpha\|\leq n$ are pairwise $1/n$-separated, because
$\langle u_\alpha,v_\alpha-v_\beta\rangle=1$. A separated subset of a
separable metric space is countable, and the union over $n$ is all of $I$.
\end{proof}

Call a multiplier direction \emph{active} if the stationary set for that
direction is nonempty and its image under $f$ contains more than one point.

\begin{lemma}[countability of active directions]\label{lem:active}
The active projective multiplier directions form a countable subset of
$\Pj$.
\end{lemma}

\begin{proof}
It is enough to work in the affine chart
$[-t:1]$, $t\in\R$; the single missing projective direction can be added at
the end. Put
\[
 B_t=A_2-tA_1,\qquad b_t=a_2-ta_1,\qquad K_t=\ker B_t.
\]
For every $t$ whose stationary set is nonempty, choose one
$x_t$ satisfying $B_tx_t+b_t=0$ and put
$g_t=A_1x_t+a_1$.

If $h\in K_t$, exact Taylor expansion gives
\begin{equation}\label{eq:fiberchange}
 f_1(x_t+h)-f_1(x_t)
 =2\langle g_t,h\rangle+\langle A_1h,h\rangle.
\end{equation}
The corresponding change in $f_2$ is $t$ times the right side, because
$f_2-tf_1$ is constant on the affine stationary set $x_t+K_t$.

First consider the set $Q$ of parameters $t$ for which the restriction of the
symmetric bilinear form $\langle A_1\cdot,\cdot\rangle$ to $K_t$ is nonzero.
For every $t\in Q$, choose $h_t,k_t\in K_t$ with
$\langle A_1h_t,k_t\rangle\ne0$. If $s\ne t$ and $k_s\in K_s$, then
\[
 t\langle A_1h_t,k_s\rangle
 =\langle A_2h_t,k_s\rangle
 =\langle h_t,A_2k_s\rangle
 =s\langle A_1h_t,k_s\rangle.
\]
Thus $\langle A_1h_t,k_s\rangle=0$. After rescaling,
$(A_1h_t,k_t)_{t\in Q}$ is a biorthogonal system, so
Lemma~\ref{lem:biorthogonal} makes $Q$ countable.

Now let $t\notin Q$. The quadratic term in \eqref{eq:fiberchange} vanishes
identically on $K_t$, so the fiber is active exactly when
$\langle g_t,\cdot\rangle|_{K_t}$ is nonzero. Let $L$ be this set of
parameters. If $s\ne t$ and $h_s\in K_s$, stationarity at $t$ gives
\[
 \langle A_2x_t+a_2,h_s\rangle
 =t\langle A_1x_t+a_1,h_s\rangle.
\]
Self-adjointness and $A_2h_s=sA_1h_s$ give
\[
 \langle A_2x_t+a_2,h_s\rangle
 =s\langle A_1x_t+a_1,h_s\rangle;
\]
here we also used
$\langle a_2-sa_1,h_s\rangle=0$, obtained by pairing
$B_sx_s+b_s=0$ with $h_s$. Hence
\[
 \langle g_t,h_s\rangle=0,\qquad s\ne t.
\]
For each $t\in L$, choose $v_t\in K_t$ with
$\langle g_t,v_t\rangle\ne0$. After rescaling,
$(g_t,v_t)_{t\in L}$ is biorthogonal. Therefore $L$ is countable.

By \eqref{eq:fiberchange}, every active parameter belongs to $Q\cup L$.
Adding the missing projective direction proves the lemma.
\end{proof}

\subsection{The area estimate}

Let $\mathcal E$ be the countable set of active directions from
Lemma~\ref{lem:active}. Fix $\eta>0$ and cover $\mathcal E$ by an open set
$O_\eta\subset P$ whose
component arcs $J_s$ have total length less than $\eta$. We may take
$\eta$ small enough that all these arcs admit aligned representatives.
The set $f(B_H(R))$ is bounded, say by $M_R$. By \eqref{eq:slab}, the values
with multiplier in $J_s$ lie in the intersection of a ball of radius $M_R$
with a strip of width $2C_R|J_s|$. Hence their planar outer measure is at
most $C'_R|J_s|$. Countable subadditivity gives
\begin{equation}\label{eq:badarcs}
 {\Leb}^{*}\bigl(f(\mathcal I_R|_{O_\eta})\bigr)\leq C'_R\eta.
\end{equation}

On the compact complement $F_\eta=P\setminus O_\eta$, the stationary point
need not be unique, but its value under $f$ is unique. The projection of
$\mathcal I_R|_{F_\eta}$ onto its compact multiplier image is a continuous
surjection from a compact space to a Hausdorff space and hence a quotient map.
Since $f$ is constant on every projection fiber, it induces a continuous
selected value
\[
 y([\lambda])=f(x_\lambda)
\]
which is uniformly continuous on the compact multiplier set where it is
defined.
Let $\omega_\eta(\delta)\to0$ be a modulus of continuity.

Partition $P$ into $O(\delta^{-1})$ arcs of length at most $\delta$. On each
arc, the selected values have diameter at most $\omega_\eta(\delta)$ and, by
\eqref{eq:slab}, normal width at most $2C_R\delta$. They therefore fit in a
rectangle of area at most
$C_R''\delta\,\omega_\eta(\delta)$. Summing over the arcs and letting
$\delta\downarrow0$ gives
\[
 \Leb\bigl(f(\mathcal I_R|_{F_\eta})\bigr)=0.
\]
Together with \eqref{eq:badarcs}, this shows
\[
 \Leb\bigl(f(\Crit(f)\cap B_H(R))\bigr)\leq C'_R\eta.
\]
Let $\eta\downarrow0$. The bounded critical-value set is null. Finally,
$\Crit(f)=\bigcup_{N\geq1}(\Crit(f)\cap B_H(N))$, so countable subadditivity
proves Theorem~\ref{thm:main}.

\section{Limitations and verification}

The source expects global Sard for all $f\in\mathscr P_2^m(H)$. The present
argument proves the expectation completely for $m=2$, including noncommuting
pencils with uncountably many singular projective directions. It makes no
claim for $m>2$. The biorthogonality step is special to a two-dimensional
target: two distinct projective multipliers force the relevant mixed pairing
to vanish, whereas in higher target dimension it need only lie in the
intersection of two hyperplanes.

The Taylor identity was independently checked by the included numerical script
on random finite-dimensional symmetric pencils. This check is not part of the
proof. The bounded literature search described in the README found no exact
statement of the unrestricted target-dimension-two theorem; novelty confidence
is moderate.

\begin{thebibliography}{9}
\bibitem{LRT}
A. Lerario, L. Rizzi, and D. Tiberio,
\emph{Sard properties for polynomial maps in infinite dimension},
arXiv:2407.02296v4, accepted in the Journal of the European Mathematical
Society.

\bibitem{Conway}
J. B. Conway,
\emph{A Course in Functional Analysis}, 2nd ed.,
Springer, 1990.
\end{thebibliography}

\end{document}
